% --- File: sections/04_experiments.tex ---

\subsection{Data Acquisition and Feature Engineering}
To empirically validate the proposed architecture, we utilize historical daily market data for Bitcoin (BTC/USD) as our extreme stress-testing proxy. As established, assets exhibiting such severe structural breaks, extreme kurtosis, and pronounced volatility clustering provide an optimal worst-case environment to validate the model's resilience against generalized macroeconomic shocks.

The raw price data is transformed into a stationary series of logarithmic returns: $r_t = \ln(P_t / P_{t-1})$. To provide the neural network with sufficient contextual state, we engineer a feature matrix $\bm{X}_t$ that includes:
\begin{itemize}
    \item Daily log returns ($r_t$).
    \item Rolling historical volatility over varying lookback periods (e.g., 7-day, 14-day, and 30-day).
    \item The Parametric VaR estimate ($\text{VaR}_{param, t}^q$), serving as the structural anchor for the custom loss function.
\end{itemize}
All features are standardized using expanding window z-score normalization to strictly prevent data leakage and look-ahead bias during the training phase.

\subsection{Walk-Forward Backtesting Protocol}
Standard randomized train/test splits are fundamentally invalid for financial time-series evaluation. To replicate a live institutional production environment, we implement a chronological walk-forward backtesting protocol. 

The models are evaluated over a chronological out-of-sample period exceeding the 250-day regulatory minimum mandated by the FRTB framework, using the full available data after the initial training window. The walk-forward process follows these steps:
\begin{enumerate}
    \item \textbf{Initialization:} The LSTM is trained on an initial historical window of $T_{train}$ days.
    \item \textbf{Forecasting:} The model predicts the 99\% VaR for day $T_{train} + 1$.
    \item \textbf{Rolling Update:} The true return for day $T_{train} + 1$ is observed. The model window shifts forward by one day, and the neural network weights are incrementally fine-tuned.
\end{enumerate}
This rolling methodology ensures that the model must constantly adapt to new market regimes, accurately penalizing "catastrophic forgetting" and static overfitting.

\subsection{Evaluation Metrics}
The competing architectures are evaluated across two intersecting dimensions: safety (regulatory compliance) and efficiency (capital optimization).

\textbf{1. Model Safety (Kupiec POF Test):} 
We utilize the Kupiec Proportion of Failures test \cite{kupiec1995techniques} to determine if the observed breach rate $\hat{p} = x / N$ (where $x$ is the number of breaches and $N$ is the total out-of-sample days) is statistically consistent with the target probability $p = 0.01$. The Likelihood Ratio (LR) test statistic is defined as:
\begin{equation}
    \text{LR}_{pof} = -2 \ln \left( \frac{p^x (1-p)^{N-x}}{\hat{p}^x (1-\hat{p})^{N-x}} \right)
\end{equation}
Under the Kupiec framework, a model passes if its observed breach rate is statistically consistent with the 1\% target; it fails if it is either significantly above (too risky) or significantly below (over-conservative, trapping excess capital).

\textbf{2. Capital Efficiency (Mean VaR):}
A model that safely passes the Kupiec test by forecasting an excessively large VaR is mathematically valid but commercially unviable. We measure capital efficiency by calculating the mean forecasted VaR over the test period. A lower absolute Mean VaR, conditional on passing the Kupiec test, represents superior capital optimization and reduced capital drag.

% TODO: In Section 5 (Results), we will inject the comparative VaR forecast plot here using:
% \begin{figure}[htbp]
%     \centering
%     \includegraphics[width=\textwidth]{compare_plot.png}
%     \caption{Rolling 99\% VaR Forecasts: Anchored LSTM vs. Historical VaR Baseline.}
%     \label{fig:var_forecast}
% \end{figure}